\documentclass[class=book, crop=false]{standalone}
% ==============================================================================
% SETUP.TEX - CONFIGURACIÓN GLOBAL DEL PROYECTO
% ==============================================================================
% Este archivo centraliza todos los paquetes y estilos.
% Debe incluirse en el preámbulo del libro principal y de los archivos standalone.
% \PassOptionsToPackage{gray}{xcolor}
% --- 1. CODIFICACIÓN E IDIOMA ---
% fix-cm habilita las fuentes Computer Modern en tamaños arbitrarios (por debajo
% de 5 pt, entre otros). Debe cargarse antes que fontenc.
\usepackage{fix-cm}
\usepackage[utf8]{inputenc}
\usepackage[T1]{fontenc}
\usepackage[spanish, es-tabla]{babel}  
\usepackage{csquotes} % Recomendado para babel y citas

\usepackage{import}
% mode=image: Usa el PDF pre-compilado, nunca intenta recompilar.
% Debes compilar cada figura manualmente antes de compilar el libro.
\usepackage[mode=image, subpreambles=false]{standalone}

% --- 2. MATEMÁTICAS Y CIENCIAS ---
\usepackage{amsmath, amssymb, amsfonts, mathtools}
\usepackage{enumitem}

% LaTeX trae \DeclareMathSizes{5}{5}{5}{5}, así que a 5 pt (\tiny) los subíndices
% salen del mismo tamaño que la base: en las etiquetas de figuras el M de
% $\omega_M$ queda desproporcionado. Se restablece la proporción habitual.
\DeclareMathSizes{5}{5}{3.7}{3}

% --- 3. DISEÑO Y GEOMETRÍA --- 
% Unificamos los márgenes para todo el documento
\usepackage[left=2.5cm,right=2.5cm,top=3cm,bottom=3cm]{geometry}
\makeatletter
\ifstandalone % Evita que geometry dañe el recorte de las figuras bajándolas 3cm
  \@ifclasswith{standalone}{crop=false}{
    % Es un capítulo/sección: Dejamos que geometry actúe.
  }{
    % Es una figura: Neutralizamos geometry para que no las recorte mal.
    \geometry{pass}
  }
\fi
\makeatother
\usepackage{parskip} % Espaciado entre párrafos sin sangría
\usepackage{float}   % Para usar [H] en figura

% Ajustes de flotantes para evitar espacios verticales exagerados
\renewcommand{\topfraction}{0.95}      % Permite que las figuras ocupen hasta el 95% superior
\renewcommand{\bottomfraction}{0.95}   % Permite que las figuras ocupen hasta el 95% inferior
\renewcommand{\textfraction}{0.05}     % Permite hojas con tan solo un 5% de texto
\renewcommand{\floatpagefraction}{0.8} % Requiere que una página de solo figuras esté 80% llena
\setcounter{topnumber}{4}
\setcounter{bottomnumber}{4}
\setcounter{totalnumber}{10}

% --- 4. GRÁFICOS Y TABLAS ---
\usepackage{graphicx}
\usepackage{booktabs} % Tablas profesionales
\usepackage{caption}  % Control de captions y \captionof
\usepackage{tikz}
\usetikzlibrary{arrows.meta, calc, angles, quotes, babel, backgrounds}
\usepackage{pgfplots}
\usepgfplotslibrary{groupplots}
\usepgfplotslibrary{external}
\usepgfplotslibrary{patchplots}
\pgfplotsset{compat=1.18}
\usepackage[american, straightvoltages]{circuitikz}

% --- 5. COLORES Y CAJAS ---

\usepackage[table]{xcolor}
\usepackage{tcolorbox}

% Definición de colores corporativos/estilo
\definecolor{utnblue}{RGB}{0, 51, 102}
\definecolor{codegreen}{rgb}{0,0.6,0}
\definecolor{codegray}{rgb}{0.5,0.5,0.5}
\definecolor{codepurple}{rgb}{0.58,0,0.82}
\definecolor{backcolour}{rgb}{0.96,0.96,0.96}

% --- 6. ENTORNOS PERSONALIZADOS (CAJAS) ---

% Caja "Resultado" (Azul Corporativo) — Definiciones, fórmulas clave, resultados importantes.
% Uso: \begin{resultado}[Fórmula de Euler] ... \end{resultado}
\newtcolorbox{resultado}[1][]{
    colback=utnblue!5!white,
    colframe=utnblue,
    title=\textbf{#1},
    fonttitle=\bfseries,
    sharp corners,
    boxrule=0.5mm
}

% Caja "Nota" (Gris) — Advertencias, aplicaciones, contexto secundario.
% Uso: \begin{nota}[Cuidado con la calculadora] ... \end{nota}
% Uso: \begin{nota} ... \end{nota}  → título por defecto "Nota"
\newtcolorbox{nota}[1][Nota]{
    colback=codegray!5!white,
    colframe=codegray,
    title=\textbf{#1},
    fonttitle=\bfseries,
    sharp corners,
    boxrule=0.5mm
}

% Caja inline para resaltar un valor y enlazarlo con su otra aparición en una tabla
% o en una cuenta: mismo color, mismo valor.
% Uso: \marcavalor{$4$}  o  \marcavalor[codegreen]{$4$}
\newtcbox{\marcavalor}[1][utnblue]{
    on line,
    boxsep=1pt, left=2pt, right=2pt, top=1pt, bottom=1pt,
    colback=#1!10!white,
    colframe=#1,
    boxrule=0.4pt,
    arc=2pt
}

% --- 7. CÓDIGO FUENTE (LISTINGS) ---
\usepackage{listings}

\lstdefinestyle{miestilo}{
    backgroundcolor=\color{backcolour},   
    commentstyle=\color{codegreen},
    keywordstyle=\color{magenta},
    numberstyle=\tiny\color{codegray},
    stringstyle=\color{codepurple},
    basicstyle=\ttfamily\footnotesize,
    breakatwhitespace=false,         
    breaklines=true,                 
    captionpos=b,                    
    keepspaces=true,                 
    numbers=left,                    
    numbersep=5pt,                  
    showspaces=false,                
    showstringspaces=false,
    showtabs=false,                  
    tabsize=2,
    frame=single,                    
    rulecolor=\color{codegray}
}

\lstset{style=miestilo}
% Soporte para tildes y ñ en bloques de código
\lstset{literate={ñ}{{\~n}}1 {á}{{\'a}}1 {é}{{\'e}}1 {í}{{\'i}}1 {ó}{{\'o}}1 {ú}{{\'u}}1}

% Operadores matemáticos personalizados

% Vector en sentido discreto (lista finita de coordenadas): negrita + flecha.
% Uso: \vect{v}, \vect{e}_k, \vect{0}.
\newcommand{\vect}[1]{\vec{\mathbf{#1}}}

% Fecha de versión y comando para capítulos standalone
\providecommand{\verdate}{\the\year.\ifnum\month<10 0\fi\the\month.\ifnum\day<10 0\fi\the\day}
\newcommand{\chapterversion}{%
    \vspace{-1.5cm}\begin{flushright}\small\textit{\color{codegray}Versión~\verdate}\end{flushright}\vspace{0.5cm}%
}

% --- 8. HIPERVÍNCULOS (DEBE SER EL ÚLTIMO PAQUETE) ---
\usepackage[colorlinks=true, linkcolor=utnblue, urlcolor=utnblue, citecolor=utnblue]{hyperref}

% --- 9. REFERENCIAS CRUZADAS ENTRE CAPÍTULOS STANDALONE ---
% Cuando se compila un capítulo standalone, xr-hyper lee los .aux de los
% otros capítulos (si existen) para resolver \ref{} a labels externos.
% En la compilación del libro completo este bloque no se activa.
\ifstandalone
  \usepackage{xr-hyper}
  \makeatletter
  \newcommand{\xrchap}[1]{%
    \edef\xrc@self{\detokenize{Capitulo#1}}%
    \edef\xrc@job{\jobname}%
    \ifx\xrc@self\xrc@job\else
      \IfFileExists{../#1capitulo/Capitulo#1.aux}%
        {\externaldocument{../#1capitulo/Capitulo#1}}{}%
    \fi
  }
  \makeatother
  \xrchap{01}\xrchap{02}\xrchap{03}\xrchap{04}
  \xrchap{05}\xrchap{06}\xrchap{07}\xrchap{08}
  \xrchap{09}\xrchap{10}\xrchap{11}\xrchap{12}
  \xrchap{13}
\fi

\begin{document}
\setcounter{chapter}{5}
\section{Convergencia y propiedades}
\label{sec:convergencia_propiedades}

La sección anterior construyó la serie de Fourier exponencial y la
aplicó al pulso rectangular. En el camino quedaron tres cuestiones
abiertas que esta sección retoma. La primera es la igualdad
$x(t) = \sum_k c_k\,e^{jk\omega_0 t}$ de la fórmula de
síntesis~\eqref{eq:sf_sintesis_box}: la obtuvimos del cálculo formal
con la base ortogonal y dejamos pendiente bajo qué hipótesis sobre
$x$ la suma infinita la reproduce punto por punto, y qué pasa en las
discontinuidades. La segunda
es el sobrepico que apareció en la reconstrucción truncada del pulso,
que al aumentar el número de armónicos se vuelve más angosto pero no
se atenúa. La tercera son las simetrías observadas en los
coeficientes del pulso ---reales para todo $k$, módulo par, fase
impar---, que anticipamos como consecuencia de la simetría conjugada
de las señales reales y de la simetría par del pulso, sin
verificarlo todavía.

La primera subsección cierra las dos primeras cuestiones, enunciando
las condiciones de Dirichlet para la convergencia y describiendo el
fenómeno de Gibbs en las discontinuidades. El resto de la sección
cataloga las propiedades que conectan operaciones sobre $x(t)$
---desplazar, rebatir, derivar, conjugar--- con operaciones sobre
los coeficientes $c_k$; entre ellas, la propiedad de conjugación
formaliza la simetría que dejó pendiente la lectura del pulso.

% ==============================================================================
% 5.3.1 Condiciones de Dirichlet y fenómeno de Gibbs
% ==============================================================================
\subsection{Condiciones de Dirichlet y fenómeno de Gibbs}
\label{subsec:dirichlet_gibbs}

La pregunta es bajo qué hipótesis sobre $x(t)$ la fórmula de
síntesis~\eqref{eq:sf_sintesis_box} reproduce efectivamente la
señal. La respuesta clásica son las llamadas \textbf{condiciones de
Dirichlet}: tres requisitos sencillos sobre $x$ que, cuando se
cumplen, garantizan que la serie converge en cada punto y que el
valor al que converge es el que uno esperaría.

\begin{resultado}[Teorema de convergencia de Dirichlet]
Sea $x(t)$ una señal periódica de período $T_0$.

\textbf{Condiciones de Dirichlet.} Supongamos que $x$:
\begin{enumerate}[label=(\roman*), leftmargin=*]
    \item es absolutamente integrable sobre un período,
          $\int_{T_0} |x(t)|\,dt < \infty$;
    \item tiene un número finito de máximos y mínimos por período;
    \item tiene un número finito de discontinuidades por período, y
          en cada una los límites laterales $x(t^-)$ y $x(t^+)$
          existen y son finitos.
\end{enumerate}

\medskip
\textbf{Convergencia.} Entonces la serie de Fourier de $x$ converge en cada $t$ al valor
\begin{equation}
    \sum_{k=-\infty}^{\infty} c_k\,e^{jk\omega_0 t}
    \;=\;
    \begin{cases}
        x(t), & \text{si $x$ es continua en $t$}, \\[4pt]
        \dfrac{x(t^-) + x(t^+)}{2}, & \text{si $x$ tiene un salto en $t$}.
    \end{cases}
    \label{eq:dirichlet}
\end{equation}
\end{resultado}

Las tres condiciones se verifican sin esfuerzo en todas las señales
periódicas que aparecen en este libro ---pulsos, ondas cuadradas y
triangulares, dientes de sierra, sinusoides rectificadas--- así que
en la práctica funcionan como un permiso casi automático para
escribir la serie. La parte interesante de~\eqref{eq:dirichlet} es la
segunda línea: en una discontinuidad, la serie no converge a
$x(t^-)$ ni a $x(t^+)$, sino al promedio aritmético de ambos límites
laterales. La serie no puede ``elegir'' un lado del salto; lo que
hace es centrarse exactamente en el medio. Es una consecuencia de
que cada término $c_k\,e^{jk\omega_0 t}$ es continuo en $t$, así que
ninguna suma finita ---ni el límite de la sucesión de sumas
finitas--- puede reproducir la discontinuidad.

Queda explicar el sobrepico observado en la reconstrucción truncada
del pulso (figura~\ref{fig:pulso_reconstruccion}). En las cercanías
de una discontinuidad de salto, la aproximación $\hat{x}_N(t)$
presenta un sobrepico cuya altura tiende a $\approx 9\%$ del tamaño
del salto y \textbf{no se atenúa al aumentar $N$}. Lo único que
cambia con $N$ es que el sobrepico se concentra en una banda más
estrecha junto al punto de salto. Este comportamiento se llama
\emph{fenómeno de Gibbs} y es compatible con Dirichlet: la
convergencia puntual no impide que falle la uniformidad en los
entornos del salto. El valor preciso del sobrepico ---un $8{,}95\%$
del salto--- sale de una integral del seno cardinal que no vamos a
desarrollar; lo que importa retener es que es una constante
universal, la misma para cualquier señal con saltos y cualquier $N$.

La figura~\ref{fig:gibbs} muestra la aproximación del pulso para
seis valores crecientes de $N$, con la altura del sobrepico medida
en cada caso. El número se mantiene cerca del 9\% en los seis
paneles; lo que cambia con $N$ es solamente el ancho de la banda
afectada.

\begin{figure}[H]
    \centering
    \begin{minipage}{\linewidth}
        \centering
        \includestandalone[width=0.85\linewidth]{figures/fig_gibbs/fig_gibbs}
        \caption{Fenómeno de Gibbs en el tren de pulsos rectangulares.
                 Cada panel muestra la aproximación truncada
                 $\hat{x}_N$ (rojo) sobre la señal original (gris) y
                 anota la altura del sobrepico medida junto al
                 flanco descendente. (a) $N=5$, (b) $N=10$,
                 (c) $N=20$, (d) $N=40$, (e) $N=80$, (f) $N=100$.
                 La altura permanece cerca del 9\% del tamaño del
                 salto en los seis casos; lo único que cambia con
                 $N$ es el ancho de la banda donde el sobrepico
                 aparece.}
        \label{fig:gibbs}
    \end{minipage}
\end{figure}

% ==============================================================================
% Bisagra: motivación y notación
% ==============================================================================

Atendida la convergencia, pasamos al catálogo de propiedades. La
motivación es doble. Por un lado, las propiedades hacen visible
cómo cada operación sobre la señal se traduce en una operación
sobre los coeficientes, dándole estructura a la correspondencia
entre tiempo y frecuencia. Por otro, son operativas: calcular los
$c_k$ desde la fórmula de análisis exige resolver la integral
parametrizada en $k$, y eso no siempre es directo; las propiedades
que siguen evitan ese trabajo cuando ya conocemos los coeficientes
de una señal y queremos los de una versión ligeramente modificada
---una combinación lineal con otra señal, un retardo, un
rebatimiento, una conjugación, una derivada o una integral---. En
lugar de reintegrar, basta aplicar la regla correspondiente a los
$c_k$ ya conocidos. Algunas propiedades además permiten leer
información directamente sobre el espectro ---retardos como
pendientes en la fase, simetrías como condiciones de paridad,
potencia repartida entre armónicos--- sin tener que volver al
dominio del tiempo.

Para enunciarlas de manera compacta usaremos la notación
\begin{equation*}
    x(t) \;\longleftrightarrow\; c_k,
\end{equation*}
con el sentido de que la sucesión $\{c_k\}_{k\in\mathbb{Z}}$ son los
coeficientes de Fourier de la señal periódica $x(t)$. La doble flecha
subraya que las fórmulas de análisis~\eqref{eq:ck_analisis_box} y de
síntesis~\eqref{eq:sf_sintesis_box} establecen una correspondencia
biunívoca: cualquiera de los dos lados determina por completo al
otro.

% ==============================================================================
% 5.3.2 Propiedad de Linealidad
% ==============================================================================
\subsection{Propiedad de linealidad}
\label{subsec:prop_linealidad}

La fórmula de análisis~\eqref{eq:ck_analisis_box} es una integral, y
las integrales son lineales. La consecuencia inmediata es que el mapa
que envía a cada señal periódica sus coeficientes de Fourier hereda
esa linealidad.

\begin{resultado}[Linealidad]
Sean $x(t)$ e $y(t)$ dos señales periódicas del mismo período $T_0$,
con coeficientes $c_k$ y $d_k$ respectivamente, y sean
$\alpha,\beta\in\mathbb{C}$. Entonces $\alpha\,x(t) + \beta\,y(t)$
es periódica con el mismo período y sus coeficientes son
\begin{equation}
    \alpha\,x(t) + \beta\,y(t)
    \;\longleftrightarrow\;
    \alpha\,c_k + \beta\,d_k.
    \label{eq:prop_linealidad}
\end{equation}
\end{resultado}

La verificación es directa desde la fórmula de análisis: al sacar la
combinación lineal afuera de la integral, los coeficientes de
$\alpha x + \beta y$ son
\begin{equation*}
    \frac{1}{T_0}\int_{T_0}\bigl[\alpha\,x(t) + \beta\,y(t)\bigr]
        e^{-jk\omega_0 t}\,dt
    \;=\; \alpha\cdot\frac{1}{T_0}\int_{T_0} x(t)\,e^{-jk\omega_0 t}\,dt
    \;+\; \beta\cdot\frac{1}{T_0}\int_{T_0} y(t)\,e^{-jk\omega_0 t}\,dt,
\end{equation*}
que son exactamente $\alpha\,c_k + \beta\,d_k$.

La condición de que los dos períodos sean iguales no es opcional. Si
$x$ tiene período $T_x$ e $y$ tiene período $T_y$ distinto, la
combinación $\alpha x + \beta y$ es periódica solo cuando el cociente
$T_x/T_y$ es racional, y su período fundamental es entonces el mínimo
común múltiplo de $T_x$ y $T_y$. La propiedad sigue valiendo en ese
caso si se reescribe a $x$ e $y$ sobre el período común; lo que no
funciona es sumar dos series con períodos distintos sin resolver
antes la cuestión del período común.

% ==============================================================================
% 5.3.3 Propiedad de Desplazamiento temporal
% ==============================================================================
\subsection{Propiedad de desplazamiento temporal}
\label{subsec:prop_desplazamiento}

Si corremos la señal $x(t)$ un tiempo $t_0$ hacia adelante, todas
sus componentes armónicas tienen que correrse lo mismo. La imagen
es mecánica: la señal es la superposición de sus armónicos, así
que retardar la señal entera es retardar cada armónico por
igual. Retardar un armónico no toca su amplitud ---sigue oscilando
con el mismo módulo, solo que llega un poco más tarde a cada
pico---, así que lo único que un desplazamiento puede modificar es
la fase. Esa es la lectura del resultado que sigue.

\begin{resultado}[Desplazamiento temporal]
Si $x(t)\leftrightarrow c_k$, entonces para cualquier $t_0\in\mathbb{R}$,
\begin{equation}
    x(t - t_0)
    \;\longleftrightarrow\;
    c_k\,e^{-jk\omega_0 t_0}.
    \label{eq:prop_desplazamiento}
\end{equation}
\end{resultado}

Para verificarlo, calculemos los coeficientes de $x(t-t_0)$ desde la
fórmula de análisis. Aplicando el cambio de variable $u = t - t_0$ ---por
lo tanto $t = u + t_0$ y $dt = du$--- la integral sobre un período en $t$
queda transformada en una integral sobre un período en $u$, gracias a
que el integrando $x(u)\,e^{-jk\omega_0(u+t_0)}$ es periódico en $u$ con
período $T_0$ (un intervalo de longitud $T_0$ es equivalente a cualquier
otro):
\begin{equation*}
    \frac{1}{T_0}\int_{T_0} x(t - t_0)\,e^{-jk\omega_0 t}\,dt
    \;=\; \frac{1}{T_0}\int_{T_0} x(u)\,e^{-jk\omega_0(u + t_0)}\,du.
\end{equation*}
El factor $e^{-jk\omega_0 t_0}$ no depende de $u$ y sale de la
integral, dejando
\begin{equation*}
    e^{-jk\omega_0 t_0}\cdot\frac{1}{T_0}\int_{T_0} x(u)\,e^{-jk\omega_0 u}\,du
    \;=\; e^{-jk\omega_0 t_0}\,c_k,
\end{equation*}
que es justamente el lado derecho de~\eqref{eq:prop_desplazamiento}.

El factor $e^{-jk\omega_0 t_0}$ tiene módulo $1$, así que
\begin{equation*}
    |c_k\,e^{-jk\omega_0 t_0}| \;=\; |c_k|,
\end{equation*}
en línea con la intuición de arriba: el espectro de módulos es
invariante por desplazamientos temporales, porque la amplitud de
cada armónico no depende de dónde elegimos el origen de tiempo.
Toda la información del retardo se acumula en la fase, y el cambio
es \emph{lineal en $k$} con pendiente $-\omega_0 t_0$: cada
armónico acumula una fase proporcional a su orden, porque a mayor
frecuencia el mismo retardo $t_0$ representa una fracción más
grande del período del armónico. Leer un retardo en el espectro es
entonces leer una pendiente en la fase.

% ==============================================================================
% 5.3.4 Propiedad de Rebatimiento
% ==============================================================================
\subsection{Propiedad de rebatimiento}
\label{subsec:prop_rebatimiento}

El rebatimiento $x(t)\mapsto x(-t)$ refleja la señal en el eje del
tiempo. La pregunta natural es qué le hace a los coeficientes; la
respuesta es la más simple posible.

\begin{resultado}[Rebatimiento]
Si $x(t)\leftrightarrow c_k$, entonces
\begin{equation}
    x(-t) \;\longleftrightarrow\; c_{-k}.
    \label{eq:prop_rebatimiento}
\end{equation}
\end{resultado}

Los coeficientes de $x(-t)$ son
\begin{equation*}
    \frac{1}{T_0}\int_{T_0} x(-t)\,e^{-jk\omega_0 t}\,dt.
\end{equation*}
Con el cambio de variable $u = -t$ se invierte el sentido de
integración, lo que cancela el signo de $du = -dt$, y el exponente
queda $e^{-jk\omega_0(-u)} = e^{j k\omega_0 u} = e^{-j(-k)\omega_0 u}$.
La integral resulta
\begin{equation*}
    \frac{1}{T_0}\int_{T_0} x(u)\,e^{-j(-k)\omega_0 u}\,du \;=\; c_{-k}.
\end{equation*}

Geométricamente, rebatir la señal en el tiempo es lo mismo que dar
vuelta el espectro: la sucesión de coeficientes se lee de derecha a
izquierda. Por sí sola la propiedad es modesta; combinada con la
conjugación que sigue da las simetrías que reducen la serie a
sumas de solo cosenos o de solo senos cuando la señal es real par o
real impar, lo que será el insumo de la sección siguiente.

% ==============================================================================
% 5.3.5 Propiedad de Conjugación
% ==============================================================================
\subsection{Propiedad de conjugación}
\label{subsec:prop_conjugacion}

La conjugación $x(t) \mapsto \overline{x(t)}$ refleja la señal en
el eje real del plano complejo. Para señales generales tiene
contenido propio. El caso que más usaremos en el libro es el de
señales reales, donde $\overline{x(t)} = x(t)$ y la propiedad se
vuelve una restricción estructural sobre los coeficientes ---la
\emph{simetría conjugada}---: es la regla que organiza todo lo que
viene en la sección de la serie trigonométrica.

\begin{resultado}[Conjugación]
Si $x(t)\leftrightarrow c_k$, entonces
\begin{equation}
    \overline{x(t)} \;\longleftrightarrow\; \overline{c_{-k}}.
    \label{eq:prop_conjugacion}
\end{equation}
\end{resultado}

Para verificarlo, conjugamos los dos lados de la fórmula de
análisis~\eqref{eq:ck_analisis_box}. El factor $1/T_0$ es real y
sale intacto del conjugado, el conjugado de la integral es la
integral del conjugado, y el conjugado de un producto es el producto
de los conjugados; aplicando estas tres observaciones,
\begin{align*}
    \overline{c_k}
    &\;=\; \overline{\,\frac{1}{T_0}\int_{T_0} x(t)\,e^{-jk\omega_0 t}\,dt\,} \\
    &\;=\; \frac{1}{T_0}\,\overline{\,\int_{T_0} x(t)\,e^{-jk\omega_0 t}\,dt\,} \\
    &\;=\; \frac{1}{T_0}\int_{T_0} \overline{x(t)\,e^{-jk\omega_0 t}}\,dt \\
    &\;=\; \frac{1}{T_0}\int_{T_0} \overline{x(t)}\,\overline{e^{-jk\omega_0 t}}\,dt.
\end{align*}
El conjugado de la exponencial cambia el signo del exponente,
$\overline{e^{-jk\omega_0 t}} = e^{jk\omega_0 t}$, que reescribimos
como $e^{-j(-k)\omega_0 t}$ para reconocer en él la exponencial de
orden $-k$ de la fórmula de análisis. Sustituyendo,
\begin{equation*}
    \overline{c_k}
    \;=\; \frac{1}{T_0}\int_{T_0}\overline{x(t)}\,e^{-j(-k)\omega_0 t}\,dt.
\end{equation*}
La integral del lado derecho es lo que la fórmula de
análisis~\eqref{eq:ck_analisis_box} produciría aplicada a la señal
$\overline{x(t)}$ con índice $-k$: es entonces el coeficiente de
orden $-k$ de $\overline{x(t)}$. La igualdad dice, por lo tanto, que
$\overline{c_k}$ es el coeficiente de orden $-k$ de $\overline{x(t)}$.
Reemplazando $k$ por $-k$, el coeficiente de orden $k$ de
$\overline{x(t)}$ es $\overline{c_{-k}}$, que es lo que afirma la
propiedad.

\paragraph{Simetría conjugada para señales reales.}
Cuando $x(t)$ es real, $\overline{x(t)} = x(t)$, así que $x(t)$ y
$\overline{x(t)}$ tienen los mismos coeficientes. La
propiedad~\eqref{eq:prop_conjugacion} dice que los coeficientes de
$\overline{x(t)}$ son $\overline{c_{-k}}$; igualándolos a los $c_k$
de $x(t)$ se obtiene $c_k = \overline{c_{-k}}$, que reemplazando
$k$ por $-k$ se escribe en la forma habitual
\begin{equation}
    c_{-k} \;=\; \overline{c_k}, \qquad k\in\mathbb{Z}.
    \label{eq:simetria_ck}
\end{equation}
Tomando módulo y argumento de los dos lados,
\begin{equation}
    |c_{-k}| \;=\; |c_k|, \qquad
    \angle c_{-k} \;=\; -\angle c_k.
    \label{eq:simetria_espectros}
\end{equation}
El espectro de módulos resulta \textbf{par} respecto del eje de
frecuencias y el espectro de fases \textbf{impar}. Toda la
información del lado positivo se replica sobre el negativo: en la
práctica suele alcanzar con dibujar la mitad $k\geq 0$ del
espectro. Ambas simetrías se ven directamente en los espectros del
pulso (figura~\ref{fig:espectro_pulso}) y del diente de sierra
(figura~\ref{fig:espectro_sierra}) calculados en la sección
anterior: los módulos son pares respecto de $k=0$ en ambos casos,
y la imparidad de la fase queda nítida en el diente de sierra, con
$+\pi/2$ a la derecha del cero y $-\pi/2$ a la izquierda. El pulso
es un caso degenerado en este punto. Como la señal es real \emph{y
además par}, sus coeficientes salen reales puros, y entonces la
fase $\angle c_k$ queda restringida a los valores $0$ o $\pi$ ---el
ángulo de un número real positivo o negativo, respectivamente---.
Estos dos ángulos tienen la particularidad de coincidir con su
opuesto: $-0 = 0$, y $-\pi$ y $\pi$ apuntan al mismo punto
(difieren en una vuelta completa). Cuando
la fase solo toma esos dos valores, las condiciones de imparidad
$\angle c_{-k} = -\angle c_k$ y de paridad $\angle c_{-k} = \angle c_k$
piden lo mismo: el lado derecho da el mismo número en ambas. Por eso graficar la fase del pulso con $\pi$
a ambos lados de $k=0$ ---como en la
figura~\ref{fig:espectro_pulso}--- satisface las dos a la vez.

\begin{resultado}[Simetría conjugada de señales reales]
Si $x(t)$ es periódica y \textbf{real}, sus coeficientes de Fourier
satisfacen
\begin{equation*}
    c_{-k} \;=\; \overline{c_k}, \qquad k\in\mathbb{Z}.
\end{equation*}
En particular, el espectro de módulos es par y el espectro de fases
es impar respecto de la frecuencia.
\end{resultado}

La lectura geométrica del por qué es directa. Si $x(t)$ es real, en
la síntesis $x(t) = \sum_k c_k\,e^{jk\omega_0 t}$ todas las
componentes con parte imaginaria deben cancelarse en pares. El par
natural es $k$ y $-k$: $c_k\,e^{jk\omega_0 t}$ es un vector giratorio
antihorario, $c_{-k}\,e^{-jk\omega_0 t}$ es uno horario, y la única
manera de que la suma de los dos caiga sobre el eje real es que
sean conjugados entre sí, es decir, $c_{-k} = \overline{c_k}$.
La figura~\ref{fig:pares_conjugados} muestra esa cancelación: dos
vectores de mismo módulo girando en sentidos opuestos, y su suma
trazando la oscilación real correspondiente. El plano complejo
aparece rotado $90^\circ$, con el eje real apuntando hacia arriba,
para alinearlo con el eje vertical del gráfico temporal: así la
altura de la punta del vector suma coincide, en cada instante,
con el valor que el panel (b) traza como función del tiempo.

\begin{figure}[H]
    \centering
    \begin{minipage}{\linewidth}
        \centering
        \includestandalone[width=0.85\linewidth]{figures/fig_pares_conjugados/fig_pares_conjugados}
        \caption{Suma de un par de vectores giratorios conjugados
                 $c\,e^{j\omega_0 t}$ y $\bar{c}\,e^{-j\omega_0 t}$:
                 imagen geométrica del par $(c_k, c_{-k})$ con
                 $c_{-k} = \overline{c_k}$ en el caso $k=1$,
                 idéntica para cualquier otro $k$.
                 (a) Instantánea en el plano complejo en
                 $t = T_0/8$, con el plano rotado $90^\circ$ para
                 que el eje real quede vertical y comparta escala
                 con el eje vertical del panel (b): los dos
                 vectores tienen el mismo módulo, giran a la misma
                 velocidad en sentidos opuestos y su suma (vector
                 negro) cae sobre el eje real. (b) Suma
                 $2|c|\cos(\omega_0 t)$ en el dominio del
                 tiempo. La línea punteada conecta la punta del
                 vector suma en (a) con el valor de la oscilación
                 en (b) en ese mismo instante; la cancelación de
                 las partes imaginarias es la imagen geométrica
                 de la simetría conjugada $c_{-k} = \overline{c_k}$.}
        \label{fig:pares_conjugados}
    \end{minipage}
\end{figure}

\paragraph{Simetrías par e impar para señales reales.}
Combinando la simetría conjugada con el rebatimiento se obtienen
las dos simetrías que aparecen de manera recurrente en los cálculos.
\begin{itemize}[leftmargin=*]
    \item \textbf{Si $x$ es real y par} ---$x(-t) = x(t)$---, la
          propiedad de rebatimiento~\eqref{eq:prop_rebatimiento} da
          $c_{-k} = c_k$, y la simetría
          conjugada~\eqref{eq:simetria_ck} da
          $c_{-k} = \overline{c_k}$. Igualando los lados derechos,
          $c_k = \overline{c_k}$: cada coeficiente coincide con su
          conjugado, lo que equivale a que sea \textbf{real para
          todo $k$}. La fase $\angle c_k$ toma entonces solo los
          valores $0$ (cuando $c_k > 0$) y $\pi$ (cuando $c_k < 0$),
          sin valores intermedios.
    \item \textbf{Si $x$ es real e impar} ---$x(-t) = -x(t)$---, el
          rebatimiento da $c_{-k} = -c_k$ y la simetría conjugada
          $c_{-k} = \overline{c_k}$, así que $c_k = -\overline{c_k}$:
          cada coeficiente es el opuesto de su conjugado, lo que
          equivale a que sea \textbf{imaginario puro para todo $k$}.
          En $k=0$ esto da $c_0 = -c_0$, y por lo tanto $c_0 = 0$
          ---consistente con que el valor medio de una señal impar
          sobre un período simétrico es nulo---. La fase $\angle c_k$
          toma solo los valores $+\pi/2$ (cuando la parte imaginaria
          es positiva) y $-\pi/2$ (cuando es negativa).
\end{itemize}

\paragraph{Volviendo a los ejemplos.}
Con la simetría conjugada en mano, podemos cerrar las observaciones
que dejamos pendientes en los ejemplos del pulso rectangular
(Subsección~\ref{subsec:ejemplo_pulso}) y del diente de sierra
(Subsección~\ref{subsec:ejemplo_sierra}). Ambas señales son reales,
así que sus coeficientes cumplen $c_{-k} = \overline{c_k}$: por eso
los dos espectros muestran módulo par en $k$ y, en el caso del
diente de sierra, fase impar nítida en $\pm\pi/2$. La diferencia
en la naturaleza de los $c_k$ viene de la simetría adicional de
cada señal. El pulso es par, así que sus coeficientes salen reales
para todo $k$, como en~\eqref{eq:ck_pulso}. El diente de sierra, al descontarle el valor medio,
queda impar; por eso sus coeficientes para $k\neq 0$ salen
imaginarios puros, como en~\eqref{eq:ck_sierra}.

% ==============================================================================
% 5.3.6 Propiedad de Diferenciación e integración
% ==============================================================================
\subsection{Propiedad de diferenciación e integración}
\label{subsec:prop_derivada_integral}

En esta subsección vemos cómo se traducen a los coeficientes la
derivada y la integral de una señal periódica. Lo notable es que
ambas operaciones se reducen a multiplicar (o dividir) cada
coeficiente por $jk\omega_0$.

\begin{resultado}[Diferenciación e integración]
Si $x(t)\leftrightarrow c_k$ y $x$ es derivable, su derivada cumple
\begin{equation}
    \dot{x}(t) \;\longleftrightarrow\; jk\omega_0\,c_k.
    \label{eq:prop_derivada}
\end{equation}
Si además $c_0 = 0$, cualquier primitiva $y(t)$ de $x(t)$ es
periódica con el mismo período y sus coeficientes para $k\neq 0$ son
\begin{equation}
    y(t) \;\longleftrightarrow\; \frac{c_k}{jk\omega_0}.
    \label{eq:prop_integral}
\end{equation}
\end{resultado}

Para~\eqref{eq:prop_derivada}, basta derivar término a término la
fórmula de síntesis~\eqref{eq:sf_sintesis_box}:
\begin{equation*}
    \dot{x}(t)
    \;=\; \frac{d}{dt}\sum_{k=-\infty}^{\infty} c_k\,e^{jk\omega_0 t}
    \;=\; \sum_{k=-\infty}^{\infty} c_k\,(jk\omega_0)\,e^{jk\omega_0 t}
    \;=\; \sum_{k=-\infty}^{\infty}\bigl(jk\omega_0\,c_k\bigr)\,e^{jk\omega_0 t}.
\end{equation*}
Comparando con la fórmula de síntesis aplicada a $\dot{x}$, se lee
que el coeficiente de orden $k$ de $\dot{x}$ es $jk\omega_0\,c_k$.

Para~\eqref{eq:prop_integral}, hay dos cuidados.

El primero es la condición $c_0 = 0$: la componente continua $c_0$
de $x(t)$ tiene primitiva $c_0\,t$, que crece linealmente y rompe
la periodicidad. Pedir $c_0 = 0$ descarta ese término y deja a la
primitiva $y(t)$ periódica. Bajo esa condición, integrando término
a término ---la primitiva de $e^{jk\omega_0 t}$ es
$e^{jk\omega_0 t}/(jk\omega_0)$, y el término $k=0$ se omite por
hipótesis---,
\begin{equation*}
    y(t)
    \;=\; \sum_{k\neq 0}\frac{c_k}{jk\omega_0}\,e^{jk\omega_0 t}
        \;+\; C.
\end{equation*}

El segundo cuidado es la constante $C$. Una primitiva está definida
salvo una constante aditiva: si $Y(t)$ es primitiva de $x(t)$,
también lo es $Y(t) + C$ para cualquier $C\in\mathbb{R}$. Esa
indeterminación se fija con una \textbf{condición inicial}: dar el
valor de $y$ en un instante, por ejemplo $y(0) = y_0$, determina
$C$ unívocamente. La constante $C$ resulta ser el valor medio de
$y(t)$ sobre un período ---es decir, su coeficiente de orden
cero---, y por eso la fórmula~\eqref{eq:prop_integral} solo cubre
los $k\neq 0$: ese coeficiente de orden cero de $y$ depende de qué
primitiva elijamos y no se puede leer de los $c_k$ de $x$.

\paragraph{Lectura.}
La derivada \emph{enfatiza} los armónicos altos: el factor
$jk\omega_0$ crece linealmente con $|k|$, así que las componentes
de alta frecuencia de $x$ pesan más en $\dot x$, que resulta una
señal más oscilante. La integración hace lo opuesto: divide por
$jk\omega_0$ y \emph{atenúa} las componentes de alta frecuencia, así que $y$
queda más suave que $x$. Operativamente, las dos propiedades
permiten leer los coeficientes de $\dot x$ y de cualquier primitiva
de $x$ a partir de los $c_k$ ---por multiplicación o división por
$jk\omega_0$---, sin tener que volver a integrar la fórmula de
análisis sobre la señal modificada.

% ==============================================================================
% 5.3.7 Propiedad — Identidad de Parseval
% ==============================================================================
\subsection{Propiedad de Parseval}
\label{subsec:prop_parseval}

Toda la maquinaria geométrica construida en la Sección~\ref{sec:producto_interno} ---el producto
interno entre señales y la ortogonalidad de la base armónica--- tiene
una consecuencia directa sobre la potencia de la señal: la potencia
media de $x(t)$ se reparte de manera limpia sobre los armónicos, sin
términos cruzados. Esa identidad es lo que se conoce como
\textbf{identidad de Parseval} para series de Fourier.

\begin{resultado}[Identidad de Parseval]
Sea $x(t)$ una señal periódica de período $T_0$ con coeficientes
$c_k$. La potencia media de $x$ es igual a la suma de las potencias
individuales de sus armónicos:
\begin{equation}
    P_x \;=\; \frac{1}{T_0}\int_{T_0} |x(t)|^2\,dt
    \;=\; \sum_{k=-\infty}^{\infty} |c_k|^2.
    \label{eq:parseval}
\end{equation}
\end{resultado}

La verificación es un cálculo en tres pasos. El primero es escribir
$|x(t)|^2 = x(t)\,\overline{x(t)}$ y reemplazar el conjugado por su
serie de síntesis, tomando conjugado en \eqref{eq:sf_sintesis_box} y
renombrando el índice de suma a $m$:
\begin{equation*}
    |x(t)|^2
    \;=\; x(t)\,\overline{x(t)}
    \;=\; x(t)\sum_{m=-\infty}^{\infty}\overline{c_m}\,e^{-jm\omega_0 t}.
\end{equation*}
El segundo paso es integrar sobre un período y sacar la suma fuera
de la integral ---intercambio válido para las señales del libro, bajo
las mismas condiciones de regularidad ya invocadas en el
Capítulo~\ref{cap04} para intercambiar integrales---:
\begin{equation*}
    \frac{1}{T_0}\int_{T_0}|x(t)|^2\,dt
    \;=\; \sum_{m}\overline{c_m}\cdot
        \frac{1}{T_0}\int_{T_0} x(t)\,e^{-jm\omega_0 t}\,dt.
\end{equation*}
El tercero es reconocer que la integral interna es exactamente el
coeficiente de orden $m$ de $x$ ---es la fórmula de análisis con $k$
reemplazado por $m$---, así que vale $c_m$. Reemplazando,
\begin{equation*}
    \frac{1}{T_0}\int_{T_0}|x(t)|^2\,dt
    \;=\; \sum_{m} \overline{c_m}\,c_m
    \;=\; \sum_{m} |c_m|^2.
\end{equation*}

La identidad de Parseval cierra la equivalencia entre las dos
representaciones de la señal: medir su potencia integrando
$|x(t)|^2$ en el tiempo o sumando $|c_k|^2$ sobre los armónicos da
exactamente el mismo número.

% ==============================================================================
% 5.3.8 Tabla resumen
% ==============================================================================
\subsection{Tabla resumen}
\label{subsec:tabla_propiedades}

Las propiedades vistas a lo largo de la sección se resumen en la
tabla~\ref{tab:propiedades_sf}. La tabla está pensada como
referencia rápida: cada fila es una regla operativa para evitar
reintegrar la fórmula de análisis cuando lo que se quiere es
relacionar los coeficientes de una señal modificada con los de la
original.

\begin{table}[H]
    \centering
    \begin{tabular}{lll}
        \toprule
        \textbf{Propiedad} & $\boldsymbol{x(t)}$ & $\boldsymbol{c_k}$ \\
        \midrule
        Linealidad        & $\alpha\,x(t) + \beta\,y(t)$
                          & $\alpha\,c_k + \beta\,d_k$ \\[4pt]
        Desplazamiento    & $x(t - t_0)$
                          & $c_k\,e^{-jk\omega_0 t_0}$ \\[4pt]
        Rebatimiento      & $x(-t)$
                          & $c_{-k}$ \\[4pt]
        Conjugación       & $\overline{x(t)}$
                          & $\overline{c_{-k}}$ \\[4pt]
        Derivación        & $\dot{x}(t)$
                          & $jk\omega_0\,c_k$ \\[4pt]
        Integración \scriptsize{($c_0=0$)}
                          & primitiva $y(t)$ de $x(t)$
                          & $\dfrac{c_k}{jk\omega_0}\;\;(k\neq 0)$ \\[8pt]
        Parseval          & $\dfrac{1}{T_0}\displaystyle\int_{T_0}|x(t)|^2\,dt$
                          & $\displaystyle\sum_{k}|c_k|^2$ \\
        \bottomrule
    \end{tabular}
    \caption{Propiedades de la serie de Fourier exponencial. En cada
             fila, $x(t)\leftrightarrow c_k$ e $y(t)\leftrightarrow d_k$
             son señales periódicas del mismo período $T_0$ con sus
             respectivos coeficientes.}
    \label{tab:propiedades_sf}
\end{table}

\end{document}
